% 06-discussion.tex — 讨论

\section{讨论}
\label{sec:discussion}

\subsection{Family Cap 的 ``均值降 / 峰值升'' Pareto 前沿}

小范围扫描的族数上限实验揭示了一条重要规律: maxFam 由 1 增至 2、4、8 时, 均值从 0.6984 单调下降至 0.6306, 但全局峰值却由 maxFam=1 的 0.7999 升至 maxFam=8 的 0.8233. 该 ``均值降 / 峰值升'' 的双目标特征与神经架构搜索 (NAS) 中的 Pareto 前沿现象一致: 更大的搜索空间提升了 ``突破单族天花板'' 的概率, 但不保证均值同步提升 \cite{liu2018darts, real2019regularized}.

从激活族数分布看, maxFam=8 的运行中 81.8\% 最终仅保留 $\text{active} \le 3$ 的族, 且全场最高分 0.8233 本身仅 active=2 (即 6/8 族被关闭). 这表明 maxFam=k 的实际作用是将 ES 初始化为 $k$ 组独立的 1 族起点, 而非真正同时激活 $k$ 族进行联合优化。当某个随机种子恰好落于 ``较优的多族组合'' 邻域时, 多族可突破单族天花板; maxFam=8 提供的是 ``获得高价值峰值突破点的能力'', 而非 ``更高的均值稳健性''.

\medskip\noindent\textbf{方法论价值.} 该发现表明: 在 Pareto 前沿上开展 ``均值-峰值双目标优化'', 显式区分 ``均值稳健'' 与 ``峰值突出'' 两类配置, 比折叠到单一标量 $maze\_quality$ 上更合理 --- 单标量难以同时表达 ``稳定可复现的局部最优'' 与 ``偶发的全局突破'' 两个目标。

\subsection{两类瓶颈机制的清晰分离}

§5.3 的消融数据揭示了两种性质截然不同的瓶颈, 该清晰分类是本文的核心方法论贡献之一:

\noindent\textbf{(1) 预算受限型 (chebyshev-4, 均值 0.157).}
chebyshev-4 邻域严格包含 chebyshev-2, 理论上 chebyshev-2 的最优解是其搜索空间的子集, 但实际 chebyshev-4 的得分低 0.59 分。原因并非 ``无解'', 而在于 chebyshev-4 的 80 位 cells mask 占据 $2^{80}$ 种取值, 叠加 birth ($2^9$) + survive ($2^9$) + priority ($2^4$) 后, 单族总搜索空间 $= 2^{103} = 2^{80} \times 2^{23}$; $40 \times 60$ 网格上 $10^5$ 的预算对其 $2^{103}$ 子空间而言过于稀疏, ES 在大量 ``无梯度平台'' 上原地踏步, 难以接近 chebyshev-2 风格的局部最优。该瓶颈源于预算约束, 是高维搜索的固有挑战, 但其 \textbf{解决方案明确}: 增加预算或降低 mask 维度即可缓解, 属于工程上可解的问题。

\noindent\textbf{(2) 表征受限型 (manhattan-1, 均值 0.421).}
manhattan-1 是 4 邻域 (von Neumann), 每个元胞每步最多感知 4 个邻居; 而走廊结构 (1-cell 宽, 长程连通) 需要 ``沿走廊方向存活 + 左右墙存活'' 的二阶条件, 4 邻域仅能 ``偶发'' 出现该模式, 难以稳定涌现。与 manhattan-2 (8 邻域, 直接编码对角线存活条件) 相差 0.358 分, 反映了直接编码对角线存活条件的能力鸿沟。该瓶颈源于模板本身的信息瓶颈, 必须更换模板才能解决, 增加预算无效。

\noindent\textbf{方法论价值.} chebyshev-4 与 manhattan-1 均得分较低, 但分属不同类别: 前者属搜索空间问题 (可通过增加预算或降低 mask 维度解决), 后者属表征能力问题 (须更换模板, 增加邻域半径或引入跳接结构). 将二者混为一谈将导致错误的调参方向: 在 manhattan-1 上增加 10 倍预算并不能将均值从 0.421 提升至 0.75, 必须将模板更换为 manhattan-2. 该清晰的二分类为后续精准调优提供了路线图 --- ``遇到低分时, 先判断属于预算问题还是表征问题, 再选择对应的解决路径''.

\subsection{大范围扫描的含义: $40 \times 60$ 网格上的性能上界}

大范围扫描 5/5 全部完成, 最高分 0.8195 与小范围扫描最高分 0.8233 几乎持平 --- 该结果具有重要意义, 揭示了 manhattan-2 / maxFam=8 配置的 \textbf{配置鲁棒性 (configuration robustness)}.

\noindent\textbf{该配置在 $40 \times 60$ 网格上的性能上界已被识别.} 将网格面积扩大 170 倍后, 峰值并未显著提升, 表明 manhattan-2 / maxFam=8 在 $40 \times 60$ 网格上已达到该配置的性能上界。5 个随机种子中, s444 (0.8195) 与 s1111 (0.8138) 落于 ``较优吸引域'' (高于均值), s5555 (0.8098) 接近均值, s3333 (0.7909) 与 s2222 (0.7758) 落于 ``较差吸引域'' (低于均值)。

5-run 均值 0.8020 与小范围扫描峰值 0.8233 相差 0.021, 处于搜索空间方差范围之内。top-2 均值 (s444 + s1111) 0.8166 与 0.8233 仅相差 0.007, 说明 ``同配置在两个尺度上达到的水平'' 接近。多吸引域结构本身也是该配置的重要发现: 它提示寻找全局最优的策略应聚焦于 ``挑选较优的种子'', 而非 ``改进搜索算法本身''.

\noindent\textbf{$2^{1648}$ 搜索空间在 $40 \times 60$ 上已被充分采样的原因.} manhattan-2 的 cells mask 为 80 位, 单族搜索空间为 $2^{103}$; maxFam=8 时全染色体 $2^{824}$ 实际有效组合 (8 active 族 $\times$ 103 bits). $10^5$ 预算相当于 $2^{17}$, 在 $2^{103}$ 上的稀疏度为 $2^{-86}$; ES 以 $200 \times 500 = 10^5$ 步能够稳定找到 0.80 附近的局部最优, 5 个随机种子的标准差 0.018 给出 ``吸引域方差'' 的一个下界。大范围扫描 5 个种子中 s444 与 s1111 命中, s3333 与 s2222 未命中, 比例与 $1/5$ 接近。

\noindent\textbf{后续方向.} 若要将分数推过 0.8233, 以下三个方向均具有明确的预期收益:
\begin{itemize}[leftmargin=*, itemsep=0pt, topsep=0pt]
  \item \textbf{更换 mask}. manhattan-3 与 chebyshev-3 的均值与 manhattan-2 接近 ($\ge 0.77$), 不同的邻域拓扑可能命中更优吸引域;
  \item \textbf{增大预算}. pop $\times$ gens 由 $10^5$ 提升至 $10^6$ (即大范围扫描的预算水平), $2^{-86}$ 的稀疏度有量级改善;
  \item \textbf{两者结合}. ``换 mask + 大范围'' 的联合空间是后续工作的开放方向, 预期收益最大.
\end{itemize}

\subsection{未来工作}

本文已建立 FamilyMask 多族 CA 编码与 $maze\_quality$ 八维几何度量的基础框架, 验证了 15-pattern 基准下的清晰判别力与两个尺度扫描下的性能上界识别, 后续工作可从以下三个方向继续推进:

\begin{itemize}[leftmargin=*, itemsep=0pt, topsep=0pt]
  \item \textbf{权重自动调优}. 当前 5 个拓扑子度量权重 $(0.20, 0.15, 0.20, 0.30, 0.15)$ 与多样性侧权重 $(0.25, 0.25, 0.50)$ 是经多轮人眼评判迭代调优的结果 (§4.2), 在视觉一致性上表现良好; 进一步可通过元学习或贝叶斯优化自动调参, 系统探索权重空间, 有望找到比当前人眼调优权重更优的局部极值。
  \item \textbf{15-pattern 扩展}. 当前基准覆盖了 6 个真迷宫生成器与 9 个常见几何反例, 已能清晰区分 mode collapse (9/9 伪迷宫取 0 分, 6/6 真迷宫均值 0.711); 可进一步加入分形迷宫、Hilbert 曲线等更复杂的拓扑反例, 验证度量的判别力上限。
  \item \textbf{服务器端 scale-up}. ES 种群规模 $200 \times 500$ 已是浏览器端 WebGPU dispatch 的合理上限, 该预算下 manhattan-2 / maxFam=8 已达到性能上界; 若部署到服务器端 GPU (pop = 500, gens = 2000), 可探索 ``更换 mask (chebyshev-3, manhattan-3) + 大预算'' 的联合空间, 将当前 0.8233 的峰值进一步上推。
\end{itemize}

FamilyMask 当前采用 ``按优先级顺序遍历'' 执行 CA 步, 等价于 ``首个触发的族决定下一态''; 该 ``早退'' 语义在结构上具有可扩展性, 后续可探索 ``多族投票''、``多族加权和'' 等推广形式, 将单族 B/S 字符串之外的规则空间纳入搜索。
